% ===== PRÉAMBULE CANONIQUE — à coller VERBATIM en tête de CHAQUE .tex =====
\documentclass[11pt,a4paper]{article}
\usepackage[utf8]{inputenc}\usepackage[T1]{fontenc}\usepackage{lmodern}
\usepackage[french]{babel}
\usepackage{amsmath,amssymb,mathtools}\usepackage{icomma}
\usepackage{siunitx}\sisetup{locale=FR}  % \qty{}{}, \si{}, \ang{}
\usepackage{xcolor}\usepackage{colortbl}
\usepackage{tikz}\usepackage{pgfplots}\pgfplotsset{compat=1.18}
\usepackage[siunitx,europeanresistors]{circuitikz} % circuits (élec.)
\usepackage[most]{tcolorbox}\usepackage{enumitem}\usepackage{array,booktabs}
\usepackage[a4paper,margin=2.2cm]{geometry}
\usetikzlibrary{arrows.meta,calc,angles,quotes,positioning,patterns,%
  backgrounds,shapes.geometric,decorations.pathmorphing,decorations.markings,intersections}
\definecolor{primary}{RGB}{222,49,99}\definecolor{primarydark}{RGB}{150,22,55}
\definecolor{defcol}{RGB}{0,114,178}\definecolor{methcol}{RGB}{52,92,148}
\definecolor{excol}{RGB}{0,135,90}\definecolor{attncol}{RGB}{200,40,55}
\tcbset{boxbase/.style={enhanced,breakable,boxrule=0.9pt,arc=2.5pt,
  left=3mm,right=3mm,top=2.3mm,bottom=2.3mm,fonttitle=\bfseries,coltitle=white}}
% --- boîtes TUTORIEL (noms EXACTS, ne pas renommer) ---
\newtcolorbox{cle}[1][]{boxbase,colback=defcol!6,colframe=defcol,
  title={\textsc{À retenir}\ifx\relax#1\relax\relax\else\textmd{ : \itshape #1}\fi}}
\newtcolorbox{methode}[1][]{boxbase,colback=methcol!7,colframe=methcol,sharp corners,
  title={\textsc{Méthode}\ifx\relax#1\relax\relax\else\textmd{ --- \itshape #1}\fi}}
\newtcolorbox{exemplebox}[1][]{boxbase,colback=excol!5,colframe=excol,
  title={\textsc{Exemple}\ifx\relax#1\relax\relax\else\textmd{ : \itshape #1}\fi}}
\newtcolorbox{piege}[1][]{boxbase,colback=attncol!6,colframe=attncol,
  title={\textsc{Piège}\ifx\relax#1\relax\relax\else\textmd{ : \itshape #1}\fi}}
\newtcolorbox{remarque}[1][]{boxbase,colback=black!4,colframe=black!45,
  title={\textsc{Remarque}\ifx\relax#1\relax\relax\else\textmd{ : \itshape #1}\fi}}
% --- boîtes EXERCICE / CORRIGÉ (noms EXACTS, auto-comptées) ---
\newtcolorbox[auto counter]{exo}[1][]{boxbase,colback=primary!4,colframe=primary,
  title={\textsc{Exercice}~\thetcbcounter\ifx\relax#1\relax\relax\else\textmd{ --- \itshape #1}\fi}}
\newtcolorbox[auto counter]{corrige}[1][]{boxbase,colback=excol!4,colframe=excol,
  title={\textsc{Corrigé}~\thetcbcounter\ifx\relax#1\relax\relax\else\textmd{ --- \itshape #1}\fi}}
\newcommand{\vv}[1]{\vec{#1}}\newcommand{\ds}{\displaystyle}
\newcommand{\R}{\mathbb{R}}\newcommand{\N}{\mathbb{N}}\newcommand{\Z}{\mathbb{Z}}
% =========================================================================
\begin{document}
\sisetup{per-mode=symbol}
% ================= CORRIGÉ FACILE =================

\begin{corrige}[énergie cinétique]
\[
E_c=\tfrac12 m v^2=\tfrac12\times 500\times(7000)^2
   =250\times\num{4.9e7}\approx\qty{1.2e10}{\joule}.
\]
\end{corrige}

\begin{corrige}[énergie potentielle]
\[
E_p=-\frac{\mathcal{G}M_T m}{r}=-\frac{\num{4.0e14}\times 500}{\num{1.0e7}}
   =-\frac{\num{2.0e17}}{\num{1.0e7}}=-\qty{2.0e10}{\joule}.
\]
Elle est \emph{négative} : le satellite est lié à la Terre.
\end{corrige}

\begin{corrige}[énergie mécanique]
$E_m=-E_c=-\qty{3.0e10}{\joule}$. Elle est négative, comme pour tout satellite lié.
\end{corrige}

\begin{corrige}[le signe des énergies]
\begin{enumerate}[label=\alph*)]
  \item $E_c>0$ (une énergie cinétique est toujours positive).
  \item $E_p<0$ (référence nulle à l'infini : un corps lié a $E_p$ négative).
  \item $E_m<0$ (car $E_m=-E_c$ : négative, le satellite est lié).
\end{enumerate}
\end{corrige}

\begin{corrige}[tout à partir de $E_c$]
\begin{enumerate}[label=\alph*)]
  \item $E_p=-2E_c=-2\times\num{1.5e10}=-\qty{3.0e10}{\joule}$.
  \item $E_m=-E_c=-\qty{1.5e10}{\joule}$ (on vérifie $E_c+E_p=\num{1.5e10}-\num{3.0e10}=-\num{1.5e10}$).
\end{enumerate}
\end{corrige}

\end{document}
